\begin{abstract}
Sự gia tăng của các cuộc tấn công mạng tinh vi đòi hỏi các Hệ thống Phát hiện Xâm nhập (IDS) không chỉ có độ chính xác cao mà còn phải cực kỳ hiệu quả để triển khai trên các thiết bị Biên (Edge) và IoT vốn hạn chế về tài nguyên. Tuy nhiên, việc phát triển một hệ thống IDS như vậy gặp phải một thách thức cơ bản: để đạt được tỷ lệ triệu hồi (recall) cao cần thiết cho việc phát hiện các cuộc tấn công ẩn nấp, chiếm tỉ lệ nhỏ trong lưu lượng mạng bị mất cân bằng trầm trọng, thường đòi hỏi các mô hình ensemble đồ sộ, tốn kém tài nguyên tính toán, điều này vi phạm các ràng buộc về độ trễ tại biên. Để khắc phục điều này, chúng tôi đề xuất một khung IDS hai giai đoạn mới đạt được cả độ chính xác phát hiện hiện đại và độ trễ suy luận cực thấp. Được đánh giá trên bộ dữ liệu chuẩn CICIDS2017, phương pháp của chúng tôi trước tiên giới thiệu một mô hình "Giáo viên" (Teacher): một tập hợp Stacking Ensemble cực kỳ mạnh mẽ kết hợp một cách chiến lược các mô hình tăng cường gradient (XGBoost, LightGBM) và các mô hình bagging (Random Forest). Được huấn luyện trên một không gian vectơ được tái cân bằng về mặt toán học thông qua kỹ thuật SMOTEENN, mô hình Teacher này trung hòa hiệu quả sự mất cân bằng dữ liệu, gần như loại bỏ các cảnh báo sai (False Positives) trong khi tối đa hóa tỷ lệ recall đối với các xâm nhập thuộc lớp thiểu số. Sau đó, để giải quyết nút thắt về tính toán, chúng tôi thực hiện quy trình Chưng cất Tri thức tối ưu hóa cho Edge (Edge-Optimized Knowledge Distillation). Chúng tôi nén các ranh giới quyết định phức tạp của mô hình ensemble Teacher cồng kềnh thành một Cây Quyết định "Học sinh" (Student) nhẹ. Kết quả thực nghiệm cho thấy mô hình Teacher đạt độ chính xác phát hiện gần như hoàn hảo ($>99,8\%$), trong khi mô hình Student sau khi được chưng cất vẫn duy trì độ chính xác cạnh tranh cao ($>99\%$) nhưng hoạt động với độ trễ suy luận giảm mạnh, giúp nó nhanh hơn nhiều cấp độ. Bằng cách thu hẹp thành công khoảng cách giữa độ chính xác của ensemble chuyên sâu và hiệu quả triển khai thực tế trên Edge, nghiên cứu này đóng góp một giải pháp phát hiện mối đe dọa thời gian thực, có khả năng mở rộng cao cho các hệ sinh thái IoT hiện đại.
\end{abstract}

